\uuid{biNN}
\exo7id{651}
\titre{exo7 651}
\auteur{gourio}
\organisation{exo7}
\datecreate{2001-09-01}
\isIndication{false}
\isCorrection{false}
\chapitre{Continuité, limite et étude de fonctions réelles}
\sousChapitre{Continuité : théorie}
\cpge{mpsi-mp2i.limites-continuite.b, pcsi.limites-continuite.b, ptsi.limites-continuite.b, mpsi-mp2i.limites-continuite.c, pcsi.limites-continuite.c, ptsi.limites-continuite.c, bcpst1.limites-continuite.c, bcpst1.limites-continuite.d, mp-mpi.topologie-espaces-normes.e, pc.espaces-vectoriels-normes.e, psi.espaces-vectoriels-normes.e, pt.fonctions-deux-trois-variables.b}
\module{Analyse}
\niveau{L1}
\difficulte{}

\contenu{
\question{
Donner un exemple de fonction continue sur $[0,1]$ non lipschitzienne, puis
de fonction continue en un seul point, puis de fonction discontinue sur les
rationnels et continue sur les irrationnels, enfin de fonction continue
telle que $f(x)\in \Rr\setminus\Qq$ si $x\in \Rr\setminus\Qq$ ou si $x=0,$ et $f(x)\in \Qq$ si $x\in
\Qq\setminus\{0\}. $ Une fonction telle que $\forall x\in \Rr,\lim\limits_{h\rightarrow
0}f(x+h)-f(x-h)=0$ est-elle continue sur $\Rr$ ? Donner un exemple de bijection
de $[0,1]$ sur $[0,1]$ discontinue en tout point.
}
}
